\bookchapter{付録A：新端末の開封「超親切」ガイド}

\booksection{開封·検品「3ステップ」：}

\begin{itemize}
  \item \textbf{付録A：新端末の開封「超親切」ガイド}1. 開封·検品「3ステップ」：
  \item \textbf{外観確認}：本体や画面に傷やへこみがないか確認。
  \item \textbf{情報照合}：箱のIMEI/シリアルナンバーと領収書のものが一致するか確認。
\end{itemize}

\booksection{アクティベーションと初期設定の「重要リマインダー」：}

\begin{itemize}
  \item \textbf{：電話をかけてイヤホンとマイクを試し、写真を撮ってカメラを確認。}2. アクティベーションと初期設定の「重要リマインダー」：
  \item \textbf{ネットワーク}：まず安定したWi-Fiに接続。
  \item \textbf{アカウント}：Apple ID、Huaweiアカウント、Xiaomiアカウントなどにログイン。
  \item \textbf{基本設定}：「デバイスを探す」をオンにし、しっかりしたロック画面のパスワードを設定。
  \item \textbf{eSIM設定}：スマートフォンがeSIMに対応している場合（iPhoneや一部のAndroidフラッグシップなど）、設定でeSIMプランを追加でき、物理SIMカードの挿入は不要です。海外渡航時はeSIMが特に便利で、現地キャリアのeSIMプランを直接購入できます。
\end{itemize}

\booksection{アプリのインストールは「セキュリティ第一」：}

\begin{itemize}
  \item \textbf{：古いスマートフォンからデータを移行する際、両方のスマートフォンのバッテリーが十分にあることを確認。iPhoneは「クイックスタート」、Androidは公式の機種変更アプリを使用。}3. アプリのインストールは「セキュリティ第一」：
  \item \textbf{公式ストアを優先}：AppleはApp Store、Androidはスマートフォン付属のアプリストアを使用。
\end{itemize}
